% BHU PYQ -- Manually transcribed & error-corrected
% Paper: BCH-214 : Cost Accounting
% Exam: B.Com. (Hons.) Semester III Examination, 2023-24

\documentclass[11pt,a4paper]{article}
\usepackage[a4paper,top=2.5cm,bottom=2.5cm,left=2.8cm,right=2.8cm,headheight=14pt]{geometry}
\usepackage{amsmath,amssymb}
\usepackage[T1]{fontenc}
\usepackage[utf8]{inputenc}
\usepackage{lmodern,microtype,fancyhdr,enumitem,hyperref,lastpage,parskip}

\hypersetup{
    colorlinks=true,
    urlcolor=black,
    linkcolor=black,
    pdftitle={BCH-214: Cost Accounting -- BHU B.Com. (Hons.) Sem III 2023-24}
}

\pagestyle{fancy}
\fancyhf{}
\renewcommand{\headrulewidth}{0.4pt}
\renewcommand{\footrulewidth}{0.4pt}
\fancyhead[L]{\small   \textit{Banaras Hindu University}}
\fancyhead[R]{\small   \textit{BCH-214: Cost Accounting}}
\fancyfoot[C]{\small Page  \thepage\ of \pageref{LastPage}}


\newlist{parts}{enumerate}{2}
\setlist[parts,1]{label=(\roman*),leftmargin=*,itemsep=5pt,topsep=3pt}
\setlist[parts,2]{label=(\alph*),leftmargin=*,itemsep=3pt,topsep=2pt}

\newcommand{\pts}[1]{\hfill   \textit{[#1]}}


\begin{document}


\begin{center}
  {\large\bfseries BANARAS HINDU UNIVERSITY}\\[2pt]
  {
\normalsize B.Com. (Hons.) --- Semester III Examination, 2023-24}\\[6pt]
  \rule{0.8\linewidth}{0.5pt}\\[6pt]
  {\large\bfseries Paper No. BCH-214: Cost Accounting}\\[4pt]
  {
\normalsize Time: 3 Hours\hfill Maximum Marks: 70}
\end{center}

\rule{\linewidth}{0.4pt}

\smallskip



\noindent   \textit{Instructions: Answer all questions. All questions carry equal marks (14 marks each).}

\bigskip




\noindent   \textbf{Question 1.}
In what essential respects is Cost Accounting different from financial accounting? How does cost accounting help in controlling costs?\pts{14}

\centerline{   \textbf{OR}}

Explain why it is necessary to apply different methods of costing for different industries and state the several methods employed?\pts{14}

\medskip\rule{\linewidth}{0.2pt}




\noindent   \textbf{Question 2.}
What is meant by overhead expense? Give various methods of allocating overheads and discuss any two of these in detail giving their merits and demerits.\pts{14}

\centerline{   \textbf{OR}}

What are the objects of Material Control? Describe briefly the various methods for pricing the issue of materials.\pts{14}

\medskip\rule{\linewidth}{0.2pt}




\noindent   \textbf{Question 3.}
The following figures are available from the financial accounts for the year ended 31st March, 2022:


\begin{center}
\small

\begin{tabular}{|l|r|}
\hline
   \textbf{Particulars} &   \textbf{Amount (Rs.)} \\
\hline
Direct Materials & 2,50,000 \\
Direct Wages & 1,00,000 \\
Factory Overheads & 3,80,000 \\
Administration Overheads & 2,50,000 \\
Selling and Distribution Overheads & 4,80,000 \\
Bad Debts & 20,000 \\
Preliminary Expenses (written off) & 7,000 \\
Legal Charges & 8,000 \\
Dividends Received & 40,000 \\
Interest on Deposits Received & 25,000 \\
Sales (1,20,000 Units) & 7,20,000 \\
Closing Stock: & \\
\quad -- Finished Stock (40,000 units) & 1,20,000 \\
\quad -- Work in Progress & 80,000 \\
\hline
\end{tabular}
\end{center}

The cost accounts reveal:

\begin{parts}
  \item Direct Material Consumption Rs. 2,80,000.
  \item Direct Wages Rs. 90,000.
  \item Factory Overheads recovered at 20\% on the prime cost.
  \item Administration Overhead @ Rs. 3 per unit of production.
  \item Selling and Distribution Overhead @ Rs. 4 per unit sold.
\end{parts}

Prepare:

\begin{parts}
  \item Costing Profit and Loss Account.
  \item Financial Profit and Loss Account.
  \item Statement Reconciling the Profits and Losses.
\end{parts}\pts{14}

\centerline{   \textbf{OR}}

What is a Cost Sheet and why is it prepared? Prepare a cost sheet with imaginary figures.\pts{14}

\medskip\rule{\linewidth}{0.2pt}

\pagebreak




\noindent   \textbf{Question 4.}
Write short notes on the following:

\begin{parts}
  \item Normal and Abnormal Waste and Scrap.
  \item Joint Products and By-products.
\end{parts}\pts{14}

\centerline{   \textbf{OR}}

Explain the working of process costing and give the salient features of the system. For which kinds of business is process costing suitable?\pts{14}

\medskip\rule{\linewidth}{0.2pt}




\noindent   \textbf{Question 5.}
What is a Contract Account? How is it prepared? Give a specimen of such an account using imaginary figures.\pts{14}

\centerline{   \textbf{OR}}

How will you value work in progress in the books of a contractor? How would this item appear in the Balance Sheet? Clearly explain through an example.\pts{14}

\vfill
\rule{\linewidth}{0.4pt}

\begin{center}\small
\href{https://bhu-exam.vercel.app/}{Click here to visit our website}\\[4pt]
\href{https://me-aryan.vercel.app/}{Click here to visit our contributor}\\[4pt]
\href{https://quashan-me.vercel.app/}{Click here to visit our contributor}
\end{center}

\end{document}
